\aufzaehlungdrei{Die Jacobi-Matrix ist

\mathdisp {\begin{pmatrix}  0   &   1  &  0  \\  -1 &  0  &  0 \\ 0  &  0 &  0  \end{pmatrix}} { . }

}{Als Normalenfeld nehmen wir $\begin{pmatrix}  x  \\ y \\ z   \end{pmatrix}$, in

\mavergleichskette
{\vergleichskette
{ P
}
{ = }{ \begin{pmatrix}  0  \\ 1 \\  0   \end{pmatrix}
}
{  }{
}
{    }{
}
{   }{
}
}
{}{}{}
bilden
\mathkor {} {\begin{pmatrix}  1  \\ 0 \\  0   \end{pmatrix}} {und} {\begin{pmatrix}  0  \\ 0\\  1   \end{pmatrix}} {}
eine Basis des Tangentialraumes. Es ist

\mavergleichskettedisp
{\vergleichskette
{ \begin{pmatrix}  0   &   1  &  0  \\  -1 &  0  &  0 \\ 0  &  0 &  0  \end{pmatrix} \begin{pmatrix}  1  \\ 0 \\  0   \end{pmatrix}
}
{ =} { \begin{pmatrix}  0  \\ -1\\  0   \end{pmatrix}
}
{ } {
}
{ } {
}
{ } {
}
}
{}{}{}
und

\mavergleichskettedisp
{\vergleichskette
{  \begin{pmatrix}  0   &   1  &  0  \\  -1 &  0  &  0 \\ 0  &  0 &  0  \end{pmatrix} \begin{pmatrix}  0  \\ 0\\  1   \end{pmatrix}
}
{ =} { \begin{pmatrix}  0  \\ 0\\  0   \end{pmatrix}
}
{ } {
}
{ } {
}
{ } {
}
}
{}{}{.}
Dabei ist
\mathl{\begin{pmatrix}  0  \\ -1\\  0   \end{pmatrix}}{} ein Vielfaches zum Normalenvektor und daher ist seine orthogonale Projektion auf den Tangentialraum ebenfalls der Nullvektor. Die Abbildung
\mathl{v \mapsto \nabla_vF}{} ist also die Nullabbildung.
}{In

\mavergleichskette
{\vergleichskette
{ Q
}
{ = }{ \begin{pmatrix}  0  \\ 0\\  1   \end{pmatrix}
}
{  }{
}
{    }{
}
{   }{
}
}
{}{}{}
bilden
\mathkor {} {\begin{pmatrix}  1  \\ 0 \\  0   \end{pmatrix}} {und} {\begin{pmatrix}  0  \\ 1 \\  0   \end{pmatrix}} {}
eine Basis des Tangentialraumes. Es ist

\mavergleichskettedisp
{\vergleichskette
{ \begin{pmatrix}  0   &   1  &  0  \\  -1 &  0  &  0 \\ 0  &  0 &  0  \end{pmatrix} \begin{pmatrix}  1  \\ 0 \\  0   \end{pmatrix}
}
{ =} { \begin{pmatrix}  0  \\ -1\\  0   \end{pmatrix}
}
{ } {
}
{ } {
}
{ } {
}
}
{}{}{}
und

\mavergleichskettedisp
{\vergleichskette
{ \begin{pmatrix}  0   &   1  &  0  \\  -1 &  0  &  0 \\ 0  &  0 &  0  \end{pmatrix} \begin{pmatrix}  0  \\ 1 \\  0   \end{pmatrix}
}
{ =} { \begin{pmatrix}  1  \\ 0 \\  0   \end{pmatrix}
}
{ } {
}
{ } {
}
{ } {
}
}
{}{}{.}
Beide Bildvektoren sind bereits tangential zu $Y$ in $Q$. Eine beschreibende Matrix zu
\mathl{v \mapsto \nabla_vF}{} ist

\mathdisp {\begin{pmatrix}  0   &   -1  \\  1  &  0 \end{pmatrix}} { . }

}

 [[Kategorie:Latexseite]]
